\documentclass[twoside,spanish,a4paper,12pt]{tfm}

% Editar la titulación
\titulacion{Máster en Inteligencia Artificial y Computación Cuántica \\ Aplicada a los Mercados Financieros, MIAX Edición 14}

% Editar el título
\title{Título del Trabajo de Fin de Máster (puede ocupar varias líneas de forma automática)}

% Si es una alumna se debe usar
% \authorlabel{Autora}
\authorlabel{Autor}
% Editar el nombre
\author{Nombre del Alumno}

% Si hay varios tutores:
% \tutorlabel{Tutores}
% \tutor{Nombre del tutor 1 \\[2mm] Nombre del tutor 2}
% Si el tutor es masculino:
\tutorlabel{Tutor}
% Si la tutora es femenina:
% \tutorlabel{Tutora}
% Editar
\tutor{Nombre del Tutor/a}

% Editar: Poner mes y año de la convocatoria de lectura del TFM
\convocatoria{Octubre 2026}

\begin{document}

% NO QUITAR ESTOS ELEMENTOS
\portada
\cleardoublepage


% Editar: Resumen en Español (obligatorio)
\begin{resumen}
  Este es el resumen del TFM. Debe ser corto (máximo media página) y cubrir los aspectos principales del modelado, la aproximación analítica y los resultados obtenidos aplicados a los mercados financieros.
\end{resumen}
\cleardoublepage

% Editar: Resumen en Inglés
\begin{abstract}
  This is the abstract of the TFM. It must be short and cover the main aspects of the application of AI and quantitative models to the financial problem.
\end{abstract}
\cleardoublepage

% Editar: Agradecimientos (opcional)
\begin{agradecimientos}
  En primer lugar, quiero agradecer a la dirección académica y a mi tutor por guiarme en este proceso.

  En segundo lugar, a la máquina de café por su inestimable contribución para que mi código llegara a compilar de madrugada.
\end{agradecimientos}
\cleardoublepage

\tableofcontents

\pagestyle{tfm}
\justify

% Las figuras se buscan en el directorio figs
% Cada capítulo está en su propio fichero tex. Ver el directorio tex.
% La bibliografía está dentro del directorio bib

\chapter{Introducción}
% Contenidos del capítulo.
% Las secciones presentadas son orientativas.

\section{Introducción}

Utiliza esta plantilla \LaTeX{} para elaborar tu memoria. 

El Apéndice~\ref{cap:apendice} recoge algunos ejemplos de edición
con \LaTeX{}.

\section{Motivación}


\section{Objetivos}


\section{Organización de la memoria}


\chapter{Datos y Definición del Problema Financiero}
% Contenidos del capítulo.
% Las secciones presentadas son orientativas.

\section{Planteamiento del problema}


\section{Problema principal}


\section{Datos utilizados}


\section{Análisis de datos}


\chapter{Estado del Arte}
% Contenidos del capítulo.
% Las secciones presentadas son orientativas.

\section{Análisis de aplicaciones similares}
% Qué aplicaciones similares hay y en qué se diferencia de ellas la propuesta


\section{Tecnologías}
% Análisis crítico de las tecnologías y sistemas de despliegue posibles y por qué se han seleccionado unas concretas.


\chapter{Metodología y Modelos Algorítmicos}
% Contenidos del capítulo.
% Las secciones presentadas son orientativas.


\chapter{Experimentación, Validación y Resultados}
% Contenidos del capítulo.
% Las secciones presentadas son orientativas.



\chapter{Conclusiones y Trabajo Futuro}
% Contenidos del capítulo.
% Las secciones presentadas son orientativas.

\section{Revisión de costes}

\section{Conclusiones}

\section{Trabajo futuro}


% Editar el fichero tex/declaracion_ia.tex
\cleardoublepage
\ia
\cleardoublepage

\pagestyle{appendix}

\appendix
\chapter{Apéndice con ejemplos}
\section{Ejemplos del lenguaje de marcado \LaTeX\ para finanzas cuantitativas}
\label{cap:apendice}

Este documento es un ejemplo de cómo usar BibTeX en la gestión de bibliografía sin perder la cordura. Aquí citamos tres elementos clave: el libro de cabecera de Machine Learning Financiero \cite{lopezdeprado}, el artículo fundacional de valoración de opciones \cite{blackscholes}, y las perspectivas de futuro cuántico \cite{quantumfinance}.
Si vas a programar los modelos, seguramente necesites el manual de Hilpisch \cite{hilpisch}\footnote{Si el modelo pierde dinero en producción, la culpa no es de la bibliografía.}.

  \textbf{Sobreajuste (Overfitting)}: el gran enemigo que acecha en el párrafo 1.

  \textit{Backtesting out-of-sample}: la prueba de fuego del párrafo 2.

  \texttt{import pandas as pd}: el inicio de todo en el párrafo 3.

  \begin{itemize}
  \item Limpieza de datos (manejo de NaNs y *outliers*).
  \item Ingeniería de características (cálculo de retornos, momentum, volatilidad).
  \end{itemize}

  % Espacio vertical
  \vspace{0.5cm}

  \begin{enumerate}
  \item Entrenamiento del modelo predictivo.
  \item Optimización de hiperparámetros.
  \item Despliegue en el entorno de trading.
  \end{enumerate}

A continuación se muestra una ecuación que asusta a los de letras pero que aquí es el pan de cada día (la ecuación diferencial parcial de Black-Scholes):

  \[ \frac{\partial V}{\partial t} + \frac{1}{2}\sigma^2 S^2 \frac{\partial^2 V}{\partial S^2} + rS \frac{\partial V}{\partial S} - rV = 0 \]

  Podemos incluir imágenes en formato: png, pdf o jpg. Si usas Python o R para generar tus gráficas, lo ideal es guardarlas directamente en pdf para no perder resolución.

  En la figura~\ref{fig:diagrama} se muestra un ejemplo de figura de ancho de página. En la figura figura~\ref{fig:diagrama2} se muestra como poner dos figuras en una misma fila.
  

  \begin{figure}[!htb]
    \centering % Siempre queda mejor centrado
    \includegraphics[width=0.8\textwidth]{logo_miax.png} % Recuerda poner esta imagen en la carpeta figs
    \caption{Esta es una figura que \LaTeX\ decide dónde colocar (floating) para que el maquetado no explote como en otros procesadores de texto.}
    \label{fig:diagrama}
  \end{figure}

  \begin{figure}[!htb]
  \begin{tabular}{cc}
    Curva ROC & Matriz de Confusión \\[2mm]
    \includegraphics[width=0.4\textwidth]{logo_miax.png} &  \includegraphics[width=0.4\textwidth]{logo_miax.png}
  \end{tabular}
    \caption{Ejmplo de dos figuras en la misma fila.}
    \label{fig:diagrama2}
\end{figure}

  \vspace*{1cm}
  Este es un ejemplo de una tabla comparativa de métricas de rendimiento:

  \begin{center}
    \begin{tabular}{|l|c|c|}
      \hline
      \textbf{Modelo Algorítmico} & \textbf{Ratio Sharpe} & \textbf{Max Drawdown} \\ \hline
      Regresión Logística & 1.20 & -15.4\% \\ \hline
      Random Forest & 1.85 & -8.2\% \\ \hline
      LSTM (Red Neuronal) & 2.10 & -6.1\% \\ \hline
    \end{tabular}
  \end{center}

\end{lstlisting}

\addcontentsline{toc}{chapter}{Bibliografía}
\bibliographystyle{unsrt}
\bibliography{bib/bibliografia}

\cleardoublepage
\contraportada

\end{document}
